% 数式サンプル: フォーマットなしで Computer Modern の数式フォントを
% \textfont に張り、TFM ベースの数式組版を確認する。cm* の TFM は
% TeX Live(kpsewhich)で解決されるため、TeX Live が必要。
%
% INITEX は数式クラス(mathcode)と数式グルー(muskip)をほとんど
% 持たないため、plain.tex 相当の最小限をここで設定する。
\catcode`\{=1 \catcode`\}=2 \catcode`\$=3 \catcode`\^=7 \catcode`\_=8
\font\tenrm=cmr10 \font\teni=cmmi10 \font\tensy=cmsy10 \font\tenex=cmex10
\font\sevenrm=cmr7 \font\seveni=cmmi7 \font\sevensy=cmsy7
\font\fiverm=cmr5 \font\fivei=cmmi5 \font\fivesy=cmsy5
\textfont0=\tenrm \scriptfont0=\sevenrm \scriptscriptfont0=\fiverm
\textfont1=\teni \scriptfont1=\seveni \scriptscriptfont1=\fivei
\textfont2=\tensy \scriptfont2=\sevensy \scriptscriptfont2=\fivesy
\textfont3=\tenex \scriptfont3=\tenex \scriptscriptfont3=\tenex
\thinmuskip=3mu
\medmuskip=4mu plus 2mu minus 4mu
\thickmuskip=5mu plus 5mu
\mathcode`\+="202B
\mathcode`\-="2200
\mathcode`\=="303D
\mathcode`\<="313C
\mathcode`\>="313E
\mathcode`\(="4028
\mathcode`\)="5029
\mathcode`\,="613B
\tenrm
\shipout\hbox{$f(x) = ax^2 + bx + c, \hskip 10pt x_1 < x_2, \hskip 10pt {a+b \over c-d}$}
\end
